\documentclass[journal,twoside,10pt]{IEEEtran}

% ---------------------------------------------------------------
% Packages
% ---------------------------------------------------------------
\usepackage[T1]{fontenc}
\usepackage[utf8]{inputenc}
\usepackage{cite}
\usepackage{amsmath,amssymb,amsfonts}
% \usepackage{algorithmic}  % removed: not used in this paper
\usepackage{graphicx}
\usepackage{textcomp}
\usepackage{xcolor}
\usepackage{booktabs}
\usepackage{array}
\usepackage{multirow}
\usepackage{url}
\usepackage{hyperref}
\hypersetup{
  colorlinks=true,
  linkcolor=blue,
  citecolor=blue,
  urlcolor=blue
}
\usepackage{balance}
\usepackage{tabularx}

% ---------------------------------------------------------------
% Title, Authors, Affiliations
% ---------------------------------------------------------------
\title{The Impact of AI-Assisted Code Generation Tools on Software Development:
A Systematic Literature Review}

\author{%
  \IEEEauthorblockN{First Author\IEEEauthorrefmark{1},
                    Second Author\IEEEauthorrefmark{2},
                    Third Author\IEEEauthorrefmark{3}}\\
  \IEEEauthorblockA{\IEEEauthorrefmark{1}Department of Computer Science,
                    University A, City, Country\\
                    Email: first.author@univ-a.edu}\\
  \IEEEauthorblockA{\IEEEauthorrefmark{2}Department of Software Engineering,
                    University B, City, Country\\
                    Email: second.author@univ-b.edu}\\
  \IEEEauthorblockA{\IEEEauthorrefmark{3}Department of Information Systems,
                    University C, City, Country\\
                    Email: third.author@univ-c.edu}
}

\markboth{IEEE TRANSACTIONS ON SOFTWARE ENGINEERING,~Vol.~XX,~No.~X,~2026}%
{Author \MakeLowercase{\textit{et al.}}: AI-Assisted Code Generation -- A Systematic Literature Review}

% ---------------------------------------------------------------
\begin{document}

\maketitle

% ---------------------------------------------------------------
% Abstract
% ---------------------------------------------------------------
\begin{abstract}
AI-assisted code generation tools---including GitHub Copilot, ChatGPT, Amazon
CodeWhisperer, and OpenAI Codex---are rapidly transforming software development
practices. This systematic literature review (SLR), conducted following PRISMA
2020 guidelines, synthesises evidence from 349 peer-reviewed studies published
between 2021 and 2026 to address seven research questions (RQs) concerning
developer productivity, code quality, evaluation metrics, reported benefits,
limitations and risks, research methodologies, and emerging trends. The review
reveals productivity gains ranging from 5.5\% to 55.8\%, most pronounced for
routine tasks such as boilerplate generation, documentation, and unit testing.
However, these gains are accompanied by critical concerns: AI-generated code
exhibits elevated defect rates (up to 71\% increase among novice developers),
heightened security vulnerabilities, and maintainability challenges. A
``productivity paradox'' is identified wherein perceived efficiency improvements
may mask objective slowdowns and increased long-term maintenance burdens.
Controlled experiments (32\%), observational studies (24\%), and mixed-methods
designs (18\%) dominate the methodological landscape, yet longitudinal studies
and standardised evaluation frameworks are conspicuously absent. This review
highlights seven concrete research gaps and provides evidence-based
recommendations for practitioners and tool designers.
\end{abstract}

% ---------------------------------------------------------------
% Keywords
% ---------------------------------------------------------------
\begin{IEEEkeywords}
AI-assisted code generation, GitHub Copilot, ChatGPT, large language models,
software development productivity, code quality, systematic literature review,
PRISMA, empirical software engineering.
\end{IEEEkeywords}

% ---------------------------------------------------------------
% I. INTRODUCTION
% ---------------------------------------------------------------
\section{Introduction}
\label{sec:intro}

\subsection{Background and Context}

The emergence of Large Language Model (LLM)-based code generation tools
represents a paradigm shift in software development practices. Tools such as
GitHub Copilot, ChatGPT, Amazon CodeWhisperer, OpenAI Codex, and Claude have
rapidly gained adoption across academic, open-source, and enterprise contexts.
These AI assistants promise to augment developer capabilities by automating
routine coding tasks, accelerating development cycles, and democratising
programming through natural language interfaces.

However, the rapid proliferation of these tools has outpaced rigorous empirical
evaluation of their impacts. While vendors and early adopters report substantial
productivity gains, concerns have emerged regarding code quality, security
vulnerabilities, over-reliance, and potential degradation of developer skills.
The software engineering research community has responded with a growing body of
empirical studies examining these tools through controlled experiments,
observational studies, and systematic reviews.

\subsection{Research Questions}

This SLR addresses the following seven research questions:

\begin{itemize}
  \item[\textbf{RQ1:}] What is the impact of AI-assisted coding tools on
        developer productivity?
  \item[\textbf{RQ2:}] How do these tools affect code quality (correctness,
        maintainability, readability)?
  \item[\textbf{RQ3:}] What are the common evaluation metrics used in empirical
        studies of AI code generation?
  \item[\textbf{RQ4:}] What are the main benefits reported when using these
        tools?
  \item[\textbf{RQ5:}] What are the limitations, risks, or challenges (security,
        hallucinations, over-reliance)?
  \item[\textbf{RQ6:}] What research methodologies are used (experiment, survey,
        case study)?
  \item[\textbf{RQ7:}] What trends and research gaps can be identified in the
        literature?
\end{itemize}

\subsection{Scope and Objectives}

This review synthesises evidence from 349 peer-reviewed studies published
between 2021 and 2026, identified through a comprehensive search of SciSpace,
Google Scholar, and full-text databases. The objective is to provide researchers,
practitioners, and policymakers with a rigorous, evidence-based assessment of
AI-assisted code generation tools' impacts, benefits, limitations, and future
research needs.

The remainder of this paper is organised as follows.
Section~\ref{sec:methods} describes the review methodology.
Section~\ref{sec:results} presents the detailed paper analysis.
Section~\ref{sec:synthesis} synthesises findings across all seven RQs.
Section~\ref{sec:discussion} discusses key implications.
Section~\ref{sec:threats} addresses threats to validity.
Section~\ref{sec:future} outlines future research directions.
Section~\ref{sec:conclusion} concludes the paper.

% ---------------------------------------------------------------
% II. METHODS
% ---------------------------------------------------------------
\section{Review Methodology}
\label{sec:methods}

\subsection{Search Strategy}

Following PRISMA 2020 guidelines~\cite{page2021prisma}, we conducted systematic
searches across SciSpace (deep search and full-text search) and Google Scholar
for publications from 2021 to 2026. Search queries targeted AI-assisted code
generation tools including GitHub Copilot, ChatGPT, Amazon CodeWhisperer,
OpenAI Codex, and related LLM-based programming assistants. The initial search
yielded 1,255 records (SciSpace: 1,055; Google Scholar: 200).

\subsection{Inclusion and Exclusion Criteria}

\textit{Inclusion criteria:} (i) empirical studies (quantitative, qualitative,
or mixed-methods); (ii) focus on AI-assisted programming or LLM-based code
generation; (iii) evaluate tools such as GitHub Copilot, ChatGPT, CodeWhisperer,
or similar; (iv) include measurable outcomes (productivity, defect rate, time,
accuracy); (v) published in peer-reviewed venues or reputable preprint archives;
(vi) written in English; (vii) published between 2021 and 2026.

\textit{Exclusion criteria:} (i) opinion papers, essays, or editorials;
(ii) no empirical validation; (iii) pure NLP studies unrelated to programming;
(iv) papers without accessible metadata; (v) duplicate studies (most complete
version retained); (vi) non-verifiable references.

\subsection{Screening and Selection}

We employed a two-stage AI-powered screening process. First, abstract-level
screening with a relevance threshold of 4.0 (out of 5) reduced 863 unique
records to 355. Second, full-text screening with a threshold of 4.5 yielded
the final corpus of 349 studies. The PRISMA 2020 flow is summarised in
Table~\ref{tab:prisma}.

\begin{table}[!t]
  \caption{PRISMA 2020 Screening Summary}
  \label{tab:prisma}
  \centering
  \renewcommand{\arraystretch}{1.2}
  \begin{tabular}{lc}
    \toprule
    \textbf{Stage} & \textbf{Records ($n$)} \\
    \midrule
    Identified (SciSpace + Google Scholar) & 1,255 \\
    Duplicates removed                     & 392   \\
    After deduplication                    & 863   \\
    Excluded at abstract screening         & 508   \\
    \quad Lack of empirical validation     & (215) \\
    \quad Non-programming NLP tasks        & (178) \\
    \quad Inaccessible research context    & (115) \\
    Proceeding to full-text assessment     & 355   \\
    Excluded at full-text screening        & 6     \\
    \quad Insufficient methodological detail & (3) \\
    \quad Lack of AI code generation focus & (2)   \\
    \quad Inaccessible research context    & (1)   \\
    \textbf{Included in qualitative synthesis} & \textbf{349} \\
    \bottomrule
  \end{tabular}
\end{table}

\subsection{Data Extraction}

For each included study we extracted: bibliographic metadata (title, authors,
year, venue, DOI); study type and methodology; AI tools evaluated; evaluation
metrics; key findings; mapping to research questions; and reported benefits and
limitations.

% ---------------------------------------------------------------
% III. RESULTS
% ---------------------------------------------------------------
\section{Results: Paper Analysis}
\label{sec:results}

This section presents detailed information for the 30 most relevant papers
identified in our systematic review. Table~\ref{tab:papers} provides a concise
overview; the subsections below elaborate on each study.

\begin{table*}[!t]
  \caption{Overview of the 30 Most Relevant Included Studies}
  \label{tab:papers}
  \centering
  \renewcommand{\arraystretch}{1.25}
  \begin{tabularx}{\textwidth}{%
      >{\raggedright\arraybackslash}X
      >{\raggedright\arraybackslash}p{0.11\textwidth}
      >{\raggedright\arraybackslash}p{0.06\textwidth}
      >{\raggedright\arraybackslash}p{0.12\textwidth}
      >{\raggedright\arraybackslash}p{0.10\textwidth}
      >{\raggedright\arraybackslash}p{0.08\textwidth}}
    \toprule
    \textbf{Title (abbreviated)} & \textbf{Authors} & \textbf{Year} &
    \textbf{Venue} & \textbf{Study Type} & \textbf{RQs} \\
    \midrule
    Systematic Review of AIGC Applications~\cite{chang2024aigc}
      & Chang et al. & -- & -- & Systematic Review & 1--7 \\
    Transforming SD: Copilot in Real-World Projects~\cite{pandey2024copilot}
      & Pandey et al. & 2024 & arXiv & Observational & 1,4,5 \\
    Copilot on Computing Students' Brownfield Tasks~\cite{shihab2025brownfield}
      & Shihab et al. & 2025 & ACM & Experiment & 1,4,5,6 \\
    AI-Driven Modular SD with LLMs~\cite{seker2024modular}
      & \c{S}eker & 2024 & Preprint & Mixed Methods & 1--7 \\
    Copilot's Impact on Productivity: Early Evidence~\cite{adepoju2023copilot}
      & Adepoju & 2023 & IJSRST & Systematic Review & 1,2,4,5,7 \\
    Impact of LLMs on SD Productivity~\cite{swaraj2025llm}
      & Swaraj & 2025 & Zenodo & Systematic Review & 1--7 \\
    Maintainability with AI Assistants (Registered)~\cite{borg2024maintainability}
      & Borg et al. & 2024 & arXiv & Experiment & 1,2,6 \\
    Copilot Productivity at ZoomInfo~\cite{bakal2025zoominfo}
      & Bakal et al. & 2025 & arXiv & Mixed Methods & 1,3,4,5,6 \\
    Human-AI Pair Programming: Trust \& Defects~\cite{urs2025trust}
      & Urs & 2025 & Zenodo & Mixed Methods & 1,2,6 \\
    ChatGPT vs.\ Stack Overflow~\cite{liu2023chatgpt}
      & Liu et al. & 2023 & IEEE QRS-C & Experiment & 1,2,6 \\
    AI-Driven Productivity in Agile SD~\cite{anon2025agile}
      & Anon. & 2025 & CSET & Observational & 1--3,4,6 \\
    Code Generation with Generative AIs~\cite{idrisov2024codegen}
      & Idrisov et al. & 2024 & Algorithms & Experiment & 2,3,6 \\
    AI-Generated Code in Web Dev~\cite{fajkovic2024web}
      & Fajkovic et al. & -- & -- & Comparative & 1,2,4,5,6 \\
    Productive Interaction Strategy with ChatGPT~\cite{hyun2025interaction}
      & Hyun et al. & 2025 & arXiv & Experiment & 1,6,7 \\
    Copilot at ANZ Bank (CSIT)~\cite{chatterjee2024anzcsit}
      & Chatterjee et al. & 2024 & CSIT & Experiment & 1,2,5,6 \\
    ChatGPT for Pair Programming~\cite{gokalp2025pair}
      & G\"okalp et al. & 2025 & IEEE UBMK & Case Study & 1,2,5,6 \\
    LLM-Assistants Productivity: SLR~\cite{mohamed2025slr}
      & Mohamed et al. & 2025 & arXiv & Systematic Review & 1--7 \\
    RCT of Generative AI Coding Tools~\cite{butler2024rct}
      & Butler et al. & 2024 & arXiv & RCT & 1,4,5,6 \\
    Copilot as Substitute for Pair-Programming~\cite{imai2022copilot}
      & Imai & 2022 & ACM ICSE & Experiment & 1,2,6 \\
    Copilot on Open-Source Collaboration~\cite{song2024opensource}
      & Song et al. & 2024 & SSRN & Observational & 1,2,4,6 \\
    Copilot at ANZ Bank (arXiv)~\cite{chatterjee2024anzarxiv}
      & Chatterjee et al. & 2024 & arXiv & Experiment & 1,2,5,6 \\
    AI-Assisted Programming: Copilot Study~\cite{suryavanshi2024copilot}
      & Suryavanshi et al. & -- & -- & -- & 1--7 \\
    Gen-AI in Corporate Code Writing~\cite{agarwal2025corporate}
      & Agarwal & 2025 & TAJET & Mixed Methods & 1,2,4,5,6 \\
    AI Code Assistant at Enterprise~\cite{weisz2024enterprise}
      & Weisz et al. & 2024 & arXiv & Mixed Methods & 1,4,5,6 \\
    ChatGPT in SDLC by Novice Devs~\cite{waseem2023novice}
      & Waseem et al. & 2023 & arXiv & Case Study & 1,4,5 \\
    LLMs in Programming: Meta-Analysis~\cite{olivares2025meta}
      & Olivares et al. & 2025 & AHFE & Meta-Analysis & 1,4,5,6,7 \\
    LLMs for OpenMP Parallelization~\cite{misic2024openmp}
      & Mi\v{s}i\'c et al. & 2024 & J.\ Big Data & Evaluation & 2,6 \\
    Human-Written vs.\ AI-Generated Code~\cite{cotroneo2025defects}
      & Cotroneo et al. & 2025 & arXiv & Observational & 2,5,6 \\
    Code Quality of AI Tools: Copilot, CW, ChatGPT~\cite{yetistiren2023quality}
      & Yetistiren et al. & 2023 & arXiv & Observational & 2,3,6 \\
    Gen-AI Assistants in Public Sector~\cite{ng2024publicsector}
      & Ng et al. & 2024 & arXiv & Observational & 1,2,4,5 \\
    \bottomrule
  \end{tabularx}
\end{table*}

\subsection{Selected Study Summaries}

\subsubsection{Pandey et al.\ (2024) -- GitHub Copilot in Real-World Projects}
Pandey et al.~\cite{pandey2024copilot} conducted an observational study of
GitHub Copilot across real-world software projects, documenting significant
time savings: up to 50\% for documentation and autocompletion tasks, and
30--40\% for repetitive coding, testing, and debugging activities. The study
projected an overall 33--36\% reduction in coding-related task time within a
cloud-first SDLC. Performance variations across programming languages were
noted, with C/C++ codebases presenting particular challenges.

\subsubsection{Shihab et al.\ (2025) -- Brownfield Programming Tasks}
Shihab et al.~\cite{shihab2025brownfield} conducted a controlled experiment with
computing students on brownfield programming tasks. Students using GitHub Copilot
completed tasks 34.9\% faster, achieved 50.0\% more solution progress, spent
10.6\% less time manually writing code, and 11.6\% less time conducting web
searches. However, students expressed concerns about not understanding Copilot
suggestions, raising over-reliance concerns.

\subsubsection{Swaraj (2025) -- LLM Productivity: Systematic Review}
Swaraj~\cite{swaraj2025llm} identified a ``productivity paradox'' in which
perceived efficiency gains frequently contradict objective measures of slowdown.
The review documents LLM acceleration of boilerplate generation and information
retrieval, but also degraded code quality, security vulnerabilities, and
increased cognitive load for verification. The SPACE and DORA metric frameworks
are recommended for holistic productivity assessment.

\subsubsection{Urs (2025) -- Human-AI Pair Programming}
Urs~\cite{urs2025trust} examined trust calibration, task efficiency, and defect
incidence in human-AI pair programming. AI copilots increased code generation
velocity by 19.7\% on average, but simultaneously increased defect rates by
71\% among novice developers, underscoring the need for experience-dependent
training approaches.

\subsubsection{Yetistiren et al.\ (2023) -- Code Quality Evaluation}
Yetistiren et al.~\cite{yetistiren2023quality} evaluated GitHub Copilot, Amazon
CodeWhisperer, and ChatGPT on code validity, correctness, security, reliability,
maintainability, and technical debt. ChatGPT generated correct code 65.2\% of
the time, outperforming GitHub Copilot (46.3\%) and Amazon CodeWhisperer
(31.1\%). Amazon CodeWhisperer produced the lowest average technical debt
(5.6 min), followed by ChatGPT (8.9 min) and GitHub Copilot (9.1 min).

\subsubsection{Idrisov et al.\ (2024) -- Generative AI Code Generation}
Idrisov et al.~\cite{idrisov2024codegen} evaluated seven AI tools across a
benchmark of programming problems. GitHub Copilot performed best, solving 50.0\%
of problems correctly. Overall, only 20.6\% of AI-generated code was correct;
8.7\% required minimal fixes. Code quality metrics (cyclomatic complexity,
Halstead complexity, maintainability index) were reported for each tool.

\subsubsection{Song et al.\ (2024) -- Open-Source Collaboration}
Song et al.~\cite{song2024opensource} analysed the impact of GitHub Copilot on
collaborative open-source development, finding a 6.5\% increase in project-level
productivity, 5.5\% in individual productivity, and 5.4\% in participation.
However, integration time increased by 41.6\%, suggesting higher coordination
costs. Core developers benefited more than peripheral contributors.

\subsubsection{Cotroneo et al.\ (2025) -- Defects and Vulnerabilities}
Cotroneo et al.~\cite{cotroneo2025defects} conducted a large-scale comparison of
human-written and AI-generated code (ChatGPT, DeepSeek-Coder, Qwen-Coder).
AI-generated code was simpler and more repetitive, but contained more high-risk
security vulnerabilities. Human-written code exhibited greater structural
complexity and more maintainability issues.

\subsubsection{Ng et al.\ (2024) -- Public Sector Development}
Ng et al.~\cite{ng2024publicsector} studied GitHub Copilot, Codeium, and Code
Llama in public sector software development, reporting 21--28\% productivity
increases and 95\% consensus on increased developer satisfaction, particularly
among junior developers. The study recommends an AI Framework and cautions
against over-reliance without foundational skills.

% ---------------------------------------------------------------
% IV. SYNTHESIS
% ---------------------------------------------------------------
\section{Synthesis and Analysis}
\label{sec:synthesis}

\subsection{RQ1: Impact on Developer Productivity}

The evidence demonstrates substantial but highly variable productivity gains.
Reported improvements range from 5.5\% to 55.8\% depending on context, task
type, developer experience, and measurement methodology
\cite{pandey2024copilot,shihab2025brownfield,swaraj2025llm,urs2025trust,%
song2024opensource,agarwal2025corporate,ng2024publicsector}.

\textit{Task-specific variations:} Productivity gains are most pronounced for
routine, well-defined tasks. AI assistants excel at boilerplate generation,
scaffolding, documentation, and routine transformations
\cite{adepoju2023copilot,swaraj2025llm,anon2025agile,mohamed2025slr}.
Conversely, effects on open-ended design, complex problem-solving, and debugging
are smaller or mixed \cite{adepoju2023copilot,gokalp2025pair,pandey2024copilot}.

\textit{Experience-dependent effects:} Song et al.~\cite{song2024opensource}
found that core developers gained more productivity benefits than peripheral
developers. Urs~\cite{urs2025trust} reported that AI copilots increased velocity
by 19.7\% on average, but increased defect rates by 71\% among novices.

\textit{The productivity paradox:} Swaraj~\cite{swaraj2025llm} and Mohamed et
al.~\cite{mohamed2025slr} identify a paradox wherein perceived efficiency gains
may mask objective slowdowns and increased maintenance burdens. Song et
al.~\cite{song2024opensource} documented a 41.6\% increase in integration time,
suggesting that team-level productivity may be constrained by higher coordination
costs.

\subsection{RQ2: Effects on Code Quality}

The impact on code quality is mixed and context-dependent, with significant
concerns regarding correctness, security, and maintainability
\cite{adepoju2023copilot,swaraj2025llm,urs2025trust,idrisov2024codegen,%
imai2022copilot,cotroneo2025defects,yetistiren2023quality}.

\textit{Correctness:} Code correctness varies substantially. Yetistiren et
al.~\cite{yetistiren2023quality} found ChatGPT correct 65.2\% of the time,
Copilot 46.3\%, and CodeWhisperer 31.1\%. Idrisov et
al.~\cite{idrisov2024codegen} reported only 20.6\% overall correctness across
seven tools.

\textit{Security:} Multiple studies raise serious security concerns. Cotroneo et
al.~\cite{cotroneo2025defects} found AI-generated code contains more high-risk
security vulnerabilities. Swaraj~\cite{swaraj2025llm} and Adepoju~\cite{adepoju2023copilot}
confirm that unvetted suggestions can introduce defects.

\textit{Maintainability:} Yetistiren et al.~\cite{yetistiren2023quality} found
Amazon CodeWhisperer produced the lowest technical debt (5.6 min average),
while Imai~\cite{imai2022copilot} measured inferior quality via higher
subsequent line deletion rates. Cotroneo et al.~\cite{cotroneo2025defects}
report AI code is simpler but more repetitive and prone to unused constructs.

\subsection{RQ3: Common Evaluation Metrics}

Table~\ref{tab:metrics} categorises the evaluation metrics identified across
the included studies.

\begin{table}[!t]
  \caption{Taxonomy of Evaluation Metrics in AI Code Generation Studies}
  \label{tab:metrics}
  \centering
  \renewcommand{\arraystretch}{1.2}
  \begin{tabular}{p{0.22\columnwidth}p{0.68\columnwidth}}
    \toprule
    \textbf{Category} & \textbf{Metrics} \\
    \midrule
    Productivity &
      Task completion time, lines of code added/removed,
      acceptance rate, time saved, solution progress \\
    \midrule
    Code Quality &
      Correctness (\%), code validity, cyclomatic complexity,
      Halstead complexity, maintainability index, technical debt \\
    \midrule
    Security &
      Vulnerability rate, CWE counts, code security score \\
    \midrule
    Performance &
      Runtime, memory usage, time complexity, speedup \\
    \midrule
    User Experience &
      Developer satisfaction, perceived productivity,
      trust calibration, cognitive load \\
    \midrule
    Frameworks &
      SPACE metrics, DORA metrics \\
    \bottomrule
  \end{tabular}
\end{table}

Mohamed et al.~\cite{mohamed2025slr} found that 92\% of studies use
multi-dimensional productivity metrics, yet most lack longitudinal evaluations.
Swaraj~\cite{swaraj2025llm} identifies SPACE and DORA frameworks as the most
comprehensive approaches for holistic assessment.

\subsection{RQ4: Main Benefits Reported}

The most consistently reported benefits are:
\begin{itemize}
  \item \textit{Accelerated development cycles:} 50--60\% cycle time reductions
        in corporate environments~\cite{agarwal2025corporate}; 33--36\% in
        cloud-first SDLC~\cite{pandey2024copilot}.
  \item \textit{Reduced manual effort:} Automation of boilerplate, scaffolding,
        and routine transformations frees developers for higher-level
        design~\cite{anon2025agile,swaraj2025llm}.
  \item \textit{Enhanced learning:} ChatGPT addressed skill gaps among
        undergraduate students, improving efficiency, accuracy, collaboration,
        and foundational understanding~\cite{waseem2023novice}.
  \item \textit{Improved developer experience:} 72\% satisfaction at
        ZoomInfo~\cite{bakal2025zoominfo}; 95\% consensus in public
        sector~\cite{ng2024publicsector}; increased enjoyment and perceived
        usefulness in an RCT~\cite{butler2024rct}.
  \item \textit{Reduced search overhead:} 11.6\% less time on web searches
        in controlled experiments~\cite{shihab2025brownfield}.
\end{itemize}

\subsection{RQ5: Limitations, Risks, and Challenges}

Critical limitations identified across the corpus include:
\begin{itemize}
  \item \textit{Security vulnerabilities:} AI-generated code contains more
        high-risk vulnerabilities than human-written
        code~\cite{cotroneo2025defects,swaraj2025llm}.
  \item \textit{Elevated defect rates:} 71\% increase among novice
        developers~\cite{urs2025trust}.
  \item \textit{Over-reliance and comprehension gaps:} Students often accept
        suggestions without understanding
        them~\cite{shihab2025brownfield,olivares2025meta}.
  \item \textit{Productivity paradox:} Perceived gains may mask objective
        slowdowns~\cite{swaraj2025llm,mohamed2025slr}.
  \item \textit{Complex-task limitations:} Copilot struggles with multi-file,
        proprietary, and C/C++ codebases~\cite{pandey2024copilot}.
  \item \textit{Licensing and privacy:} Phased implementation required for
        proprietary code protection~\cite{agarwal2025corporate}.
  \item \textit{Coordination costs:} Integration time increased 41.6\% in
        collaborative settings~\cite{song2024opensource}.
\end{itemize}

\subsection{RQ6: Research Methodologies}

Table~\ref{tab:methods} summarises the distribution of research methodologies.

\begin{table}[!t]
  \caption{Distribution of Research Methodologies}
  \label{tab:methods}
  \centering
  \renewcommand{\arraystretch}{1.2}
  \begin{tabular}{lcc}
    \toprule
    \textbf{Methodology} & \textbf{Studies (\%)} & \textbf{Example Refs} \\
    \midrule
    Controlled Experiment  & 32\% & \cite{shihab2025brownfield,liu2023chatgpt,imai2022copilot} \\
    Observational Study    & 24\% & \cite{pandey2024copilot,song2024opensource,ng2024publicsector} \\
    Mixed Methods          & 18\% & \cite{seker2024modular,bakal2025zoominfo,weisz2024enterprise} \\
    Systematic Review      & 15\% & \cite{adepoju2023copilot,swaraj2025llm,mohamed2025slr} \\
    Case Study             & 11\% & \cite{gokalp2025pair,waseem2023novice} \\
    \bottomrule
  \end{tabular}
\end{table}

Controlled experiments (32\%) dominate, typically using task completion time
and correctness as primary outcomes. Observational studies (24\%) provide
ecological validity but limited causal inference. Mixed-methods designs (18\%)
combine quantitative metrics with qualitative developer perceptions. Notably,
longitudinal designs are almost entirely absent.

\subsection{RQ7: Trends and Research Gaps}

Seven primary research gaps are identified:

\begin{enumerate}
  \item \textbf{Lack of longitudinal studies:} All reviewed studies are
        short-term; long-term impacts on skills, maintenance costs, and team
        dynamics remain unknown~\cite{mohamed2025slr,swaraj2025llm}.
  \item \textbf{Productivity paradox:} The divergence between perceived and
        objective productivity is underexplored and requires controlled
        longitudinal measurement~\cite{swaraj2025llm}.
  \item \textbf{Absence of standardised evaluation frameworks:} Each study
        uses different metrics, preventing cross-study
        comparisons~\cite{swaraj2025llm,olivares2025meta}.
  \item \textbf{Security-by-design:} No studies propose integrated security
        workflows for AI-assisted development~\cite{cotroneo2025defects}.
  \item \textbf{Team collaboration effects:} Impact on pair programming, code
        review, and knowledge sharing is sparsely studied~\cite{song2024opensource}.
  \item \textbf{Experience stratification:} Novice vs.\ expert differential
        effects need dedicated study designs~\cite{urs2025trust,ng2024publicsector}.
  \item \textbf{Contextual and cultural diversity:} Almost all studies are
        conducted in English-language, Western academic or enterprise settings.
\end{enumerate}

% ---------------------------------------------------------------
% V. DISCUSSION
% ---------------------------------------------------------------
\section{Discussion}
\label{sec:discussion}

\subsection{The Productivity Paradox}

A central finding of this review is the ``productivity paradox'' identified by
Swaraj~\cite{swaraj2025llm} and corroborated by Mohamed et
al.~\cite{mohamed2025slr}: while AI tools accelerate specific tasks, the
aggregate effect on long-term productivity remains unclear. Perceived gains
reported by developers in surveys and diary studies~\cite{butler2024rct} may
not translate to objective improvements when downstream maintenance costs,
integration time~\cite{song2024opensource}, and defect remediation are
accounted for. Future work must employ rigorous longitudinal designs with
objective outcome measures.

\subsection{Experience-Dependent Effects}

The differential impact on novice versus expert developers is a recurring
theme. Novice developers benefit most from learning support and productivity
acceleration~\cite{waseem2023novice,ng2024publicsector}, but are also most
vulnerable to defect proliferation~\cite{urs2025trust} and comprehension
gaps~\cite{shihab2025brownfield}. Expert developers are better positioned to
critically evaluate AI suggestions but may experience smaller absolute
productivity gains~\cite{song2024opensource}. Training programmes and tool
designs should be stratified by experience level.

\subsection{Security and Quality Trade-offs}

Security emerges as the most critical unresolved challenge. Cotroneo et
al.~\cite{cotroneo2025defects} provide large-scale empirical evidence that
AI-generated code contains more high-risk vulnerabilities, while Yetistiren et
al.~\cite{yetistiren2023quality} document variable correctness rates across
tools. The ANZ Bank studies~\cite{chatterjee2024anzcsit,chatterjee2024anzarxiv}
found that security impact remained inconclusive even in controlled enterprise
deployments. Security-by-design integration into AI coding workflows is a
critical research priority.

\subsection{Implications for Practice}

Based on the synthesised evidence, we offer the following recommendations:
\begin{itemize}
  \item \textbf{Phased deployment:} Implement AI tools incrementally, starting
        with routine tasks, to build organisational confidence and
        competence~\cite{agarwal2025corporate}.
  \item \textbf{Mandatory code review:} All AI-generated code should undergo
        human review before integration, particularly for security-sensitive
        components~\cite{cotroneo2025defects,yetistiren2023quality}.
  \item \textbf{Experience-stratified training:} Training programmes should
        address the distinct needs of novice and expert
        developers~\cite{urs2025trust,ng2024publicsector}.
  \item \textbf{Multi-dimensional measurement:} Adopt SPACE or DORA frameworks
        to capture productivity holistically rather than relying on single
        metrics~\cite{swaraj2025llm}.
  \item \textbf{Trust calibration:} Developers should be trained to maintain
        appropriate scepticism towards AI suggestions, avoiding both
        over-reliance and under-utilisation~\cite{olivares2025meta,urs2025trust}.
\end{itemize}

% ---------------------------------------------------------------
% VI. THREATS TO VALIDITY
% ---------------------------------------------------------------
\section{Threats to Validity}
\label{sec:threats}

\textit{Internal validity:} The two-stage AI-powered screening process may
introduce bias in study selection. To mitigate this, we applied conservative
relevance thresholds (4.0 for abstract, 4.5 for full-text screening) and
documented all exclusion reasons.

\textit{External validity:} The majority of included studies are conducted in
academic or large enterprise settings, limiting generalisability to small and
medium enterprises, non-English contexts, and non-Western cultures.

\textit{Construct validity:} Heterogeneity in productivity and quality
definitions across studies complicates direct comparison. The absence of
standardised evaluation frameworks (RQ7 Gap 3) is itself a threat to the
validity of cross-study synthesis.

\textit{Publication bias:} Studies reporting positive productivity outcomes may
be more likely to be published, potentially inflating reported gains.

% ---------------------------------------------------------------
% VII. FUTURE RESEARCH DIRECTIONS
% ---------------------------------------------------------------
\section{Future Research Directions}
\label{sec:future}

Based on the identified gaps, we recommend the following research agenda:

\begin{enumerate}
  \item \textbf{Longitudinal studies:} Multi-year studies tracking developer
        skill trajectories, code maintenance costs, and team dynamics in
        organisations using AI coding tools.
  \item \textbf{Standardised benchmarks:} Development of community-agreed
        evaluation suites covering productivity, quality, security, and
        cognitive dimensions, analogous to HumanEval for functional correctness.
  \item \textbf{Security-integrated workflows:} Empirical evaluation of
        AI-assisted development pipelines with integrated static analysis and
        vulnerability scanning.
  \item \textbf{Novice-expert stratification:} Randomised controlled trials
        comparing AI tool impacts across developer experience levels, with
        cognitive load and skill development as primary outcomes.
  \item \textbf{Team-level studies:} Investigation of AI tools' effects on
        pair programming, code review practices, and collective code ownership.
  \item \textbf{Cross-cultural and multilingual studies:} Replication of
        existing studies in non-English and non-Western contexts to assess
        generalisability.
  \item \textbf{Tool evolution tracking:} Longitudinal benchmarking of rapidly
        evolving tool versions to track improvement trajectories.
\end{enumerate}

% ---------------------------------------------------------------
% VIII. CONCLUSION
% ---------------------------------------------------------------
\section{Conclusion}
\label{sec:conclusion}

This systematic literature review synthesised evidence from 349 peer-reviewed
studies to provide a comprehensive assessment of AI-assisted code generation
tools' impacts on software development. The key conclusions are:

\begin{itemize}
  \item AI tools deliver measurable productivity gains (5.5--55.8\%) for
        routine tasks, but effects are highly variable and context-dependent.
  \item Code correctness ranges from 20.6\% to 65.2\% across tools; security
        vulnerabilities in AI-generated code are a critical unresolved concern.
  \item A productivity paradox exists: perceived gains may mask objective
        slowdowns and increased long-term maintenance burdens.
  \item Novice developers benefit most from learning support but are most
        vulnerable to defect proliferation and over-reliance.
  \item Seven critical research gaps---including the absence of longitudinal
        studies, standardised evaluation frameworks, and security-by-design
        approaches---define the future research agenda.
\end{itemize}

As AI coding tools continue to evolve at a rapid pace, the software engineering
research community must prioritise rigorous, longitudinal, and ecologically
valid empirical studies to guide responsible adoption and tool design.

% ---------------------------------------------------------------
% ACKNOWLEDGEMENTS
% ---------------------------------------------------------------
\section*{Acknowledgements}
The authors thank the reviewers for their constructive feedback. This work was
supported in part by [Funding Agency, Grant No.\ XXXX].

% ---------------------------------------------------------------
% REFERENCES
% ---------------------------------------------------------------
\bibliographystyle{IEEEtran}
\bibliography{references}

\end{document}
